\chapter*{结语}
\addcontentsline{toc}{chapter}{结语}

多重网格真正统一的对象不是某一种 V-cycle 伪代码，而是一条从数值语义到机器执行的连续链。数学上，smoother 负责压低局部/高频误差，粗空间负责表示它留下的慢模态；实现上，这些操作变成状态版本、dependency DAG、数据所有权、物理放置与调度；软件上，它们又落实为 hierarchy construction、level state、cycle driver、transfer、bottom solve 与 ownership redistribution。

因此性能问题也不能只用“线程越多越快”或“GPU 峰值更高”解释。Fine levels 往往暴露 bandwidth 与 locality，coarse levels 则逐渐暴露 task scarcity、launch/barrier、message latency、global reduction 和 active-resource 过剩。真正的优化对象是
\begin{equation}
\boxed{
\text{convergence mechanism}
\rightarrow
\text{data movement}
\rightarrow
\text{dependency/coordination}
\rightarrow
\text{ownership topology}
\rightarrow
\text{time-to-solution}.
}
\end{equation}

本书的软件部分进一步给出一个可复用的源码阅读原则：先问谁构造 hierarchy、谁拥有 level state、谁驱动 cycle，再追具体 kernel。AMReX MLMG、hypre BoomerAMG 与 PETSc PCMG/GAMG 虽然数据结构和控制流不同，但都可以映射回同一组数学对象。这样，profiling 才不会停在一个“大函数占比”，而能继续定位到 smoother、residual、transfer、communication、redistribution 或 bottom solve。

Benchmark 与诊断则负责约束结论的证据强度。受控实现实验用于隔离后端，平台最优实验用于工程选择；连续 MMS 与离散 MMS 验证不同层面的正确性；高 wait time、低 occupancy 或高 cache miss 都只是症状，只有竞争假设经过单变量实验排除以后才能称为根因。

最后，求解器架构不能脱离整个应用。粒子、场、反应、几何更新等 phase 可能拥有不同的自然并行粒度；一个局部最快的 Poisson solver 如果要求昂贵的 ownership conversion，也可能让完整 application step 变慢。成熟的并行多重网格设计因此不是选择某个“最佳框架”，而是在数值收敛、数据生命周期、资源拓扑与软件维护之间建立一套可验证、可复用、可随机器和问题变化而重新求解的工程模型。
